\documentclass[UTF8,11pt]{ctexart}
\usepackage[a4paper,margin=24mm]{geometry}
\usepackage{amsmath,amssymb,amsthm,mathtools,mathrsfs}
\usepackage[all]{xy}
\usepackage{xurl,hyperref,enumitem}
\usepackage{tikz}
\usetikzlibrary{arrows.meta,cd,decorations.markings,calc,patterns}
\hypersetup{hidelinks,breaklinks=true}
\setlength{\emergencystretch}{3em}
\allowdisplaybreaks
\newtheorem*{theorem}{Theorem}
\newtheorem*{thm}{Theorem}
\newtheorem*{lemma}{Lemma}
\newtheorem*{proposition}{Proposition}
\newtheorem*{prop}{Proposition}
\newtheorem*{corollary}{Corollary}
\newtheorem*{cor}{Corollary}
\newtheorem*{claim}{Claim}
\theoremstyle{definition}
\newtheorem*{definition}{Definition}
\newtheorem*{defn}{Definition}
\newtheorem*{remark}{Remark}
\newtheorem*{example}{Example}
\newcommand{\R}{\mathbb R}
\newcommand{\C}{\mathbb C}
\newcommand{\Z}{\mathbb Z}
\newcommand{\Q}{\mathbb Q}
\newcommand{\N}{\mathbb N}
\newcommand{\F}{\mathbb F}
\newcommand{\D}{\mathbb D}
\newcommand{\bB}{\mathbb B}
\newcommand{\bC}{\mathbb C}
\newcommand{\bR}{\mathbb R}
\newcommand{\bZ}{\mathbb Z}
\newcommand{\bN}{\mathbb N}
\newcommand{\bQ}{\mathbb Q}
\newcommand{\bD}{\mathbb D}
\newcommand{\bF}{\mathbb F}
\newcommand{\CP}{\mathbb{CP}}
\newcommand{\RP}{\mathbb{RP}}
\newcommand{\sO}{\mathscr O}
\newcommand{\sI}{\mathscr I}
\newcommand{\sF}{\mathscr F}
\newcommand{\sG}{\mathscr G}
\newcommand{\sH}{\mathscr H}
\newcommand{\sR}{\mathscr R}
\newcommand{\sM}{\mathscr M}
\newcommand{\sV}{\mathscr V}
\newcommand{\sU}{\mathscr U}
\DeclareMathOperator{\Hom}{Hom}
\DeclareMathOperator{\Ext}{Ext}
\DeclareMathOperator{\Tor}{Tor}
\DeclareMathOperator{\Spec}{Spec}
\DeclareMathOperator{\Proj}{Proj}
\DeclareMathOperator{\supp}{supp}
\DeclareMathOperator{\im}{im}
\DeclareMathOperator{\id}{id}
\DeclareMathOperator{\Tr}{Tr}
\DeclareMathOperator{\tr}{tr}
\DeclareMathOperator{\rank}{rank}
\DeclareMathOperator{\ord}{ord}
\DeclareMathOperator{\dist}{dist}
\DeclareMathOperator{\End}{End}
\DeclareMathOperator{\Aut}{Aut}
\DeclareMathOperator{\Gal}{Gal}
\DeclareMathOperator{\vol}{vol}
\DeclareMathOperator{\Cl}{Cl}
\DeclareMathOperator{\coker}{coker}
\DeclareMathOperator{\codim}{codim}
\DeclareMathOperator{\divergence}{div}
\DeclareMathOperator{\Ric}{Ric}
\DeclareMathOperator{\grad}{grad}
\DeclareMathOperator{\Hess}{Hess}
\newcommand{\abs}[1]{\left\lvert#1\right\rvert}
\newcommand{\sabs}[1]{\lvert#1\rvert}
\newcommand{\babs}[1]{\bigl\lvert#1\bigr\rvert}
\newcommand{\Babs}[1]{\Bigl\lvert#1\Bigr\rvert}
\newcommand{\bbabs}[1]{\biggl\lvert#1\biggr\rvert}
\newcommand{\BBabs}[1]{\Biggl\lvert#1\Biggr\rvert}
\newcommand{\norm}[1]{\left\lVert#1\right\rVert}
\newcommand{\snorm}[1]{\lVert#1\rVert}
\newcommand{\bnorm}[1]{\bigl\lVert#1\bigr\rVert}
\newcommand{\Bnorm}[1]{\Bigl\lVert#1\Bigr\rVert}
\newcommand{\bbnorm}[1]{\biggl\lVert#1\biggr\rVert}
\newcommand{\BBnorm}[1]{\Biggl\lVert#1\Biggr\rVert}
\newcommand{\widebar}[1]{\overline{#1}}
\newcommand{\nicefrac}[2]{\frac{#1}{#2}}
\newcommand{\linnprod}[2]{\langle#1,#2\rangle}
\newcommand{\myindex}[1]{#1}
\newcommand{\glsadd}[1]{}
\newcommand{\avoidbreak}{}
\newcommand{\sourceref}[1]{\textnormal{[原文：\nolinkurl{#1}]}}
\newcommand{\thmref}[1]{\sourceref{#1}}
\newcommand{\lemmaref}[1]{\sourceref{#1}}
\newcommand{\propref}[1]{\sourceref{#1}}
\newcommand{\corref}[1]{\sourceref{#1}}
\newcommand{\defnref}[1]{\sourceref{#1}}
\newcommand{\exampleref}[1]{\sourceref{#1}}
\newcommand{\exerciseref}[1]{\sourceref{#1}}
\newcommand{\figureref}[1]{\sourceref{#1}}
\newcommand{\sectionref}[1]{\sourceref{#1}}
\newcommand{\appendixref}[1]{\sourceref{#1}}
\newcommand{\stacksref}[2]{\href{https://stacks.math.columbia.edu/tag/#1}{#2 [#1]}}

\begin{document}
\section*{043\quad 辛纤维丛总空间的Thurston辛形式构造}
\subsection*{来源与计数目标}
Kartik Venkatram 与 Denis Auroux 合作整理的 MIT 18.966 \emph{Geometry of Manifolds} 人类课程讲义，Spring 2007。课程页面明确记载该署名关系；不是将学生笔记冒充授课者独著。Lecture 21，\S2，Definition 1、Theorem 1 及其跨页完整证明，pp.~1--2。获取日期：2026-09-28。
\par\url{https://ocw.mit.edu/courses/18-966-geometry-of-manifolds-spring-2007/0f8a98aeced5e4bbfc83c5a03be2fdd7_lect21.pdf}
\par\url{https://ocw.mit.edu/courses/18-966-geometry-of-manifolds-spring-2007/pages/lecture-notes/}
\par\textbf{本文件只计一个问题：}由纤维辛形式的上同调延拓类，在总空间上构造属于 $c+kf^*[\omega_B]$ 的辛形式；不附加 Lefschetz 奇点版本。
\subsection*{作者命题与人类证明}
\paragraph{2. Symplectic Fibrations}
Let $f:M\to B$ be a locally trivial fibration, with generic fiber $(F,\omega_F)$ a symplectic manifold.
\begin{definition}[1]
$f$ is symplectic if the structure group reduces to $\operatorname{Symp}(F,\omega_F)$, i.e. $\exists$ local trivializations $f:f^{-1}(U_i)\cong U_i\times F\to U_i$ s.t., over $U_i\cap U_j$, the change of trivialization is a symplectomorphism.
\end{definition}
Now, let $f:M\to B$ be a compact, locally trivial symplectic fibration with symplectic fiber $(F,\omega_F)$ and symplectic base $(B,\omega_B)$.
\begin{theorem}[1 (Thurston)]
If $\exists c\in H^2(M,\R)$ s.t. $c|_F=[\omega_F]=H^2(F,\R)$. Then $\forall k\gg0$, $\exists$ a symplectic form on $M$ in the class $c+k.f^*[\omega_B]$ for which the fibers of $f$ are symplectic submanifolds.
\end{theorem}
\begin{proof}
Choose a closed 2-form $\eta$ on $M$ s.t. $[\eta]=c$, and a cover $\{U_i\}$ of $B$ by contractible subsets with trivializations $\phi_i:f^{-1}(U_i)\to F\times U_i$ s.t. $\phi_i\circ\phi_j^{-1}$ are symplectomorphisms over $U_i\cap U_j$. Let $p_i=\operatorname{pr}_2\circ\phi_i:f^{-1}(U_i)\to F$. Then, on $U_i\times F$, $\eta$ and $p_i^*\omega_F$ are closed 2-forms, and
\[
[\eta|_{f^{-1}(U_i)}]=[p_i^*\omega_F]\in H^2(f^{-1}(U_i),\R)\cong H^2(F,\R) \tag{3}
\]
since $c|_F=[\omega_F]$. Thus, $\exists$ a 1-form $\alpha_i$ on $f^{-1}(U_i)$ s.t. $p_i^*\omega_F=\eta+d\alpha_i$. Now, let $\rho_i$ be a partition of unity on $B$ by smooth functions $\rho_i:B\to[0,1]$, $\operatorname{Supp}(\rho_i)\subset U_i$, and set $\widetilde\eta=\eta+\sum d((\rho_i\circ f)\alpha_i)$. Then $\widetilde\eta$ is closed, with $[\widetilde\eta]=[\eta]=c$: moreover,
\begin{align*}
\widetilde\eta|_{F_p=f^{-1}(p)}
&=\eta|_{F_p}+\sum_i(\rho_i\circ f)d\alpha_i|_{F_p}\\
&=\sum_i\rho_i(f(p))(\eta|_{F_p}+d\alpha_i|_{F_p})=\omega_F \tag{4}
\end{align*}
in the trivializations $\phi_i$.

We have obtained a closed 2-form $\widetilde\eta$ on $M$ s.t. $[\widetilde\eta]=c$ which is symplectic on the fibers. $\forall x\in M$, split $T_xM=V_x\oplus H_x$, where $V_x=\ker df_x$ is the tangent space to the fiber and $H_x=\{v\in T_xM\mid\widetilde\eta(v,v')=0\ \forall v'\in V_x\}$. These two spaces are in direct sum since $\widetilde\eta|_{V_x}$ is nondegenerate. $f^*\omega_B$ is nondegenerate on $H_x$ because $df_x:H_x\xrightarrow{\sim}T_{f(x)}B$, so $\widetilde\eta+kf^*\omega_B$ is nondegenerate on $H_x$ for $k\gg0$ since nondegeneracy is an open condition (consider $f^*\omega_B+\frac1k\widetilde\eta$). It is also nondegenerate on $V_x$ since $(\widetilde\eta+kf^*\omega_B)|_{V_x}=\widetilde\eta|_{V_x}$. Thus, we obtain our desired symplectic form on $M$.
\end{proof}
\subsection*{整理说明（非作者正文）}
原正文第1页的三个记号不一致按图像照录：题句写 $[\omega_F]=H^2(F,\R)$；$F\times U_i$ 的纤维投影写作 $\operatorname{pr}_2$；式(4) 的基点系数写作 $\rho_i(f(p))$。本题使用范围由前后正文明确为 $c$ 限制到纤维等于 $[\omega_F]$、$p_i$ 是到 $F$ 的投影、纤维上分割统一系数为基点值；此说明不冒充作者原句，不对三个字样作静默替换。第1页图像已取得，第2页实际取得解析正文而截图失败，未声称逐页图像核对。

标准先修为 de Rham 上同调、可缩底上的丛平凡化、分割统一、辛正交补及紧致性下非退化性的稳定性。具体修正 $\eta+\sum d((\rho_i\circ f)\alpha_i)$、保持上同调类、水平与垂直方向的非退化性均为本题正文承担的工作。难度依据是把上同调延拓条件转化为全局几何结构，而非直接套用现成存在定理。后续 Gompf、Donaldson 定理在原页没有完整证明，不计入。
\subsection*{可修改的简短标注（非作者正文）}
$c$：辛纤维丛总空间的辛结构存在性

$a$：纤维限制的上同调类；局部恰当形式；分割统一；辛正交水平子空间；大参数非退化构造

\texttt{cross\_theory\_status = yes}。上同调类相等使局部纤维辛形式之差成为恰当形式，经分割统一拼成全局闭形式，再用基空间形式控制水平退化并返回总空间辛结构。位置：Theorem 1 全证，pp.~1--2。

全部标注均为非穷尽、可修改初稿；未标注不表示未使用。
\subsection*{许可与修改记录}
材料来自 MIT OpenCourseWare，作者与课程见来源。按该站所示 CC BY-NC-SA 4.0 许可复用，整理部分沿用同一许可。\url{https://creativecommons.org/licenses/by-nc-sa/4.0/}。2026-09-28：助手转录题证、调整数学排版并加入来源、范围和初稿标注；不暗示作者或 MIT 认可本次整理。所留为本次题证转录摘录，不是原 PDF 下载件。
\end{document}
